\emph{(i)} By Theorem \claimref{singularity.modular-structure-completed-pair}, $\sigma_{k} = 2N_{k} - 1$, so $\mathcal{D}$ is generated by the commuting projectors $N_{k}$ and is spanned by the $2^{n}$ joint-number projectors $P_{\nu}$, which are pairwise orthogonal, sum to $1$ and are linearly independent; hence $\dim\mathcal{D} = 2^{n}$, and $\mathcal{D}$ is spanned equally by the monomials $\sigma_{S}$. Each $E_{k}$ is the descent of Proposition \claimref{singularity.descent-what-it-erases} for the mode $k$, applied in the window; since the $N_{k}$ commute, $E_{1}\circ\cdots\circ E_{n}(x) = \sum_{\nu}P_{\nu}xP_{\nu}$, by expanding the products. This map fixes $\mathcal{D}$ (each $P_{\nu}$ commutes with $\mathcal{D}$ and $\sum_{\nu}P_{\nu} = 1$), its range lies in the commutant of all $P_{\nu}$, which is $\mathcal{D}$ (in the Fock representation the $P_{\nu}$ are the rank-one projections onto the joint-number basis), it is idempotent and $\mathcal{D}$-bimodular, hence a conditional expectation onto $\mathcal{D}$. The script builds three modes in $\Cl(3,3)$, checks $2N_{k} = 1 + \sigma_{k}$, that $\mathcal{E}$ is idempotent on all $64$ basis blades with range of rank $8$, that it fixes the $\sigma$ monomials and the eight $P_{\nu}$ (idempotent, summing to $1$), and that it kills $v_{0}^{*}v_{1}$ and $v_{0}^{*}v_{1}^{*}$.

\emph{(ii)} If $H \in \mathcal{D}$ then $H$ commutes with each $P_{\nu}$, so $\mathcal{E}(Hx - xH) = \sum_{\nu}(HP_{\nu}xP_{\nu} - P_{\nu}xP_{\nu}H) = H\mathcal{E}(x) - \mathcal{E}(x)H$. Conversely suppose $\mathcal{E}\circ\operatorname{ad}_{H} = \operatorname{ad}_{H}\circ\mathcal{E}$ and apply both sides to $P_{\mu}$. Since $\mathcal{E}(P_{\mu}) = P_{\mu}$ the right side is $[H, P_{\mu}]$; the left side is $\sum_{\nu}P_{\nu}(HP_{\mu} - P_{\mu}H)P_{\nu} = P_{\mu}HP_{\mu} - P_{\mu}HP_{\mu} = 0$. So $[H, P_{\mu}] = 0$ for every $\mu$, and $H$ lies in the commutant of the $P_{\mu}$, which is $\mathcal{D}$. For the group: if $H \in \mathcal{D}$ then $e^{itH} \in \mathcal{D}$ and $\operatorname{Ad}(e^{itH})$ commutes with $\mathcal{E}$ by the same computation; conversely, differentiating $\alpha_{t}\circ\mathcal{E} = \mathcal{E}\circ\alpha_{t}$ at $t = 0$ in finite dimension gives $\operatorname{ad}_{H}\circ\mathcal{E} = \mathcal{E}\circ\operatorname{ad}_{H}$. The words for all $t$ matter: at a single time covariance is weaker, since a unitary permuting the joint-number basis up to phases commutes with $\mathcal{E}$; the script checks that the hop rotation $U(\theta) = \exp(\theta(v_{0}^{*}v_{1} - v_{1}^{*}v_{0}))$ commutes with $\mathcal{E}$ at $\theta = \pi/2$, where it exchanges the two one-particle states, and not at $(\cos\theta, \sin\theta) = (\tfrac35, \tfrac45)$. The script verifies the equivalence without the argument: it forms the linear map $H \mapsto \mathcal{E}\circ\operatorname{ad}_{H} - \operatorname{ad}_{H}\circ\mathcal{E}$ on the $64$-dimensional algebra of three modes, evaluated on all $64$ basis elements, finds its rank to be $64 - 8$, so its kernel has dimension $8$, and checks that the eight independent $\sigma$ monomials lie in it; the kernel is therefore exactly $\mathcal{D}$.

\emph{(iii)} Premise \claimref{singularity.l-boundary-mode-locality} states that the internal dynamics is generated by the $\sigma_{k}$ and their products, that is by elements of $\mathcal{D}$, and Proposition \claimref{singularity.number-diagonal-dynamics-models-scrambling-derived} records that this makes every generator diagonal in the joint number basis. By (ii) this is equivalent to (T) on every finite window. For a hop or a pairing term $H$ on modes $0, 1$, the script checks $\mathcal{E}(\operatorname{ad}_{H}N_{0}) = 0$ while $\operatorname{ad}_{H}N_{0} \neq 0$, and that the odd term $v_{2} + v_{2}^{*}$ fails at $N_{2}$; a diagonal generator with cross terms, $3\sigma_{0}\sigma_{1} - \tfrac27\sigma_{2} + \sigma_{0}\sigma_{1}\sigma_{2}$, passes on all $64$ basis elements. The general witness is the computation in (ii): for $H \notin \mathcal{D}$ some $[H, P_{\mu}] \neq 0$ while $\mathcal{E}([H, P_{\mu}]) = 0$.

\emph{(iv)} In the basis $\ket{n_{h}n_{\mathrm{in}}}$ the beam splitter sends $\ket{00} \mapsto \ket{00}$ and $\ket{10} \mapsto c\ket{10} + s\ket{01}$, so for a state $\varrho$ of the mode $h$ with the in-mode empty, tracing out the second slot gives $\Phi(\varrho)$ as displayed. Then $E(\Phi(\varrho)) = \operatorname{diag}(a + s^{2}d, c^{2}d) = \Phi(E(\varrho))$ and the occupation is $c^{2}d$. The script checks $\Phi\circ E = E\circ\Phi$ on all four matrix units for $(c, s) = (\tfrac35, \tfrac45)$, the occupation map $N \mapsto \tfrac{9}{25}N$ and trace preservation on a sample state, and that $\Phi$ scales the coherence by $c$ rather than erasing it, so covariance does not make the channel the descent.
